\documentclass[10pt,a4paper]{amsart}
\usepackage[margin=19mm]{geometry}
\usepackage{amsmath,amssymb,mathtools}
\usepackage[scr=rsfs,cal=euler]{mathalfa}
\usepackage{xfrac}
\usepackage[unicode,hidelinks,bookmarks=false]{hyperref}
\newcommand{\ie}{i.e.\ }
\newcommand{\cf}{cf.\ }
\newcommand{\resp}{respectively\ }
\newcommand{\Subsection}{\subsection}
\providecommand{\nobreakdash}{\nobreak\hbox{-}}
\setlength{\emergencystretch}{2em}
\multlinegap0pt
\usepackage{bm}
\usepackage{bbm}
\usepackage{dsfont}
\usepackage{xspace}
\usepackage{cancel}
\usepackage{mathtools}
\usepackage{amsthm}
\makeatletter
\@ifundefined{theorem}{\newtheorem{theorem}{Theorem}}{}
\@ifundefined{lemma}{\newtheorem{lemma}[theorem]{Lemma}}{}
\@ifundefined{proposition}{\newtheorem{proposition}[theorem]{Proposition}}{}
\makeatother
\def\one{{\mathbf{1}}}
\newcommand{\E}{\mathbb{E}}
\newcommand{\Ord}{\operatorname{Ord}}
\newcommand{\Res}{\operatorname{Res}}
\newcommand{\R}{\mathbb{R}}
\renewcommand{\d}{\, \text{d}}
\title[Trimmed strong laws and distributional limits for exponentia]{Trimmed strong laws and distributional limits for exponentially mixing systems}
\author{Max Auer, Sixu Liu}
\date{Source: arXiv:2601.08126v2 (CC BY 4.0)}
\begin{document}
\maketitle

\noindent\emph{Source and licence.} Trimmed strong laws and distributional limits for exponentially mixing systems. Version arXiv:2601.08126v2 (CC BY 4.0), distributed under the Creative Commons Attribution 4.0 International licence, as recorded in the arXiv licence metadata for this exact version and shown on the abstract page https://arxiv.org/abs/2601.08126v2. Full rights notice: licenses/arxiv-2601.08126-CC-BY-4.0.txt.\par\smallskip

\noindent\textit{Editorial scope.} Counted unit: the Proposition whose source label is \texttt{momentprop} in the article (source0.tex lines 821--838, source label \texttt{momentprop}), delivered together with the article's own complete original proof of it (source0.tex lines 858--1261, one \texttt{proof} environment of 404 lines). No mathematics is written, repaired or re-derived: statement and proof are the article's own text line for line. The proof is detached from the statement in the source: the article states the result at lines 821--838 and prints its proof at lines 858--1261, and the proof environment's own header names the statement, so the association is the article's own. Required context retained uncounted: regasm2 (lines 675--677); regadaptmeasure (lines 686--691); slowrecanlem (lines 698--703); truncvarlem (lines 721--726). Every definition, display and label of those lines is retained unchanged. Omitted from this package and not redistributed: 1--674, 678--685, 692--697, 704--720, 727--820, 839--857, 1262--3173. Nothing inside a retained excerpt is omitted or altered: every retained body is delivered exactly as the article's lines are written, so no omission receipt of any kind accompanies this package. Citations: the retained excerpts carry no citation command at all, so no bibliography entry of the article is transcribed and no reference is redistributed. Reference adaptation: none was needed and none was made; every reference inside the delivered unit resolves inside it, so no unresolved reference or citation is produced. Printed slips kept literal: no printed word, symbol, hypothesis or number of the counted statement or of its proof was altered. Layout and class: the article's own class is \texttt{amsart} with packages including inputenc, amsmath, amssymb, amsthm, babel, tikz-cd, url, enumerate, ulem, parskip, color, xcolor, hyperref, bm; the wrapper is the lane's pinned \texttt{amsart} editorial wrapper at 10pt on A4, which restates the theorem-like environments the retained text needs and adds the article's own macro definitions that the retained text needs (\texttt{\textbackslash E}, \texttt{\textbackslash Ord}, \texttt{\textbackslash R}, \texttt{\textbackslash Res}, \texttt{\textbackslash d}, \texttt{\textbackslash one}), with extra wrapper packages (bm, bbm, dsfont, xspace, cancel, mathtools). The delivered numbering is the unit's own. Assets: no retained span contains a figure environment, a graphics-inclusion command, a PDF-page inclusion, a table or a TikZ picture; no image file is copied into this package, the package declares no asset dependency and no image file is redistributed with it. Licence: version arXiv:2601.08126v2 of this article is distributed under the Creative Commons Attribution 4.0 International licence, as recorded in the arXiv licence metadata for this exact version and shown on the abstract page https://arxiv.org/abs/2601.08126v2. A full rights notice is delivered at licenses/arxiv-2601.08126-CC-BY-4.0.txt. Characters: every retained excerpt is pure ASCII, so the delivered engine, XeLaTeX with its default Latin Modern text font, reports no missing character.
\begin{equation}\label{regasm2}
\|\one_{A_N}-h_N\|_{L^1} \ll \epsilon_N r_N^a, \quad \text{and} \quad \|h_N\|_{C^{\kappa}} \ll \epsilon_N^{-\kappa} r_N^{-a'}
\end{equation}

\begin{lemma}\label{regadaptmeasure}
Let $A_N\subset M$ be a sequence of sets that is regularly adapted. Then
$$
\mu(A_N)\sim B_d\rho(x^*)r_N^d.
$$
\end{lemma}

\begin{lemma}\label{slowrecanlem}
Let $A_N\subset M$ be a sequence of sets that is regularly adapted. Then, for all constants $A,\,C>0$, there exists $N_0>1$ such that for all $N>N_0$ and all positive integers $1\le n\le C|\log r_N|$, one has
\begin{equation*}
\mu \left(A_N \cap T^{-n}(A_N) \right) \leq r_N^d |\log r_N|^{-A}.
\end{equation*}
\end{lemma}

\begin{lemma}[Covariance for the generalized truncation]\label{truncvarlem}
Let $\beta>\frac{d}{2}$, let $A_N \subset M$ be regularly adapted, and set $f_N=f \cdot \left(1-\one_{A_N}\right).$ Then for any constants $A,\,C>0$, there exists $N_0>1$ such that for all $N>N_0$ and positive integers $1\leq n\leq C |\log r_N|$,
\begin{equation}\label{truncvarlemconc}
\E\left(f_N \cdot f_N\circ T^n\right) \leq r_N^{- 2\beta + d} |\log r_N|^{-A}.
\end{equation}
\end{lemma}

\begin{proposition}\label{momentprop}
Fix $m_1,m_2\in \mathbb N$, and set $m=m_1+m_2.$ Assume that $m\geq2$ and that the system is exponentially mixing of order $m$.

Suppose that $\Ord_{x^*}(f)=\beta>\frac{d}{2}$. Let $A_N\subset M$ be a sequence of sets that is regularly adapted and satisfies $N r_N^{d} \rightarrow \infty$. Define $f_N=f \cdot \left(1-\one_{A_N}\right).$ Then
\begin{equation}\label{momentclaim}
\lim_{N\rightarrow\infty} \frac{\int_M \left(S_N(\one_{A_N}) - N \mu(A_N) \right)^{m_1} \left(S_N(f_N) - N \E( f_N)\right)^{m_2}\d\mu}{\Res_{x^*}(f)^{m_2}\bigl(B_d \rho(x^*)\bigr)^{\frac{m}{2}} N^{\frac{m}{2}} r_N^{\frac{md}{2}-m_2\beta }}=\omega(m_1,m_2),
\end{equation}
where $\alpha=\frac{\beta}{d}$,
\begin{equation*}
\sigma^2= \frac{1}{2\alpha-1},
\end{equation*}
and \begin{equation*}
\omega(m_1,m_2)= \begin{cases}
(m_1 - 1)!! (m_2-1)!! \sigma^{m_2} & \text{if both } m_1 \text{ and } m_2 \text{ are even},\\
0 & \text{otherwise}.
\end{cases}
\end{equation*}
\end{proposition}

\begin{proof}[Proof of Proposition \ref{momentprop}]
Denote $g_1=\one_{A_N}$ and $g_2=f_N$. By replacing $f$ with $f/\Res_{x^*}(f)$, we may assume throughout the proof that $\Res_{x^*}(f)=1$. The general case is recovered by rescaling at the end.






Step 1: Smooth approximation.

Let $\epsilon_N>0$ be chosen later. Using \eqref{regasm2}, there is a $C^{\kappa}$ smooth function $h_1$ such that $\one_{A_N}\le h_1\le 1$, with
\begin{equation*}
\|g_1- h_1\|_{L^1} \ll \epsilon_N r_N^a \quad \text{and} \quad \|h_1\|_{C^{\kappa}} \ll \epsilon_N^{-\kappa} r_N^{-a'}.
\end{equation*}
Define $h_2= f \cdot (1-h_1)$. Then
\begin{equation*}
\|g_2 - h_2 \|_{L^1} \ll \epsilon_N r_N^{a-\beta} \quad \text{and} \quad \|h_2\|_{C^{\kappa}} \ll \epsilon_N^{-\kappa} r_N^{-a'-\beta-\kappa}.
\end{equation*}


Define $\bar{g}_j=g_j-\int_M g_j \d\mu$ and $\bar{h}_j=h_j - \int_M h_j \d\mu$ for $j=1, 2.$ We next justify that replacing $g_j$ with $h_j$ in the integral of $\prod_{j=1}^2 \Bigl(S_N(\bar{g}_j)\Bigr)^{m_j}$ introduces a negligible error.

Since centering preserves the $L^1$-order, we have
\begin{equation}\label{cL1e}
\|\bar g_1-\bar h_1\|_{L^1}\ll \epsilon_N r_N^a
\quad \text{and} \quad
\|\bar g_2-\bar h_2\|_{L^1}\ll \epsilon_N r_N^{a-\beta}.
\end{equation}


We now estimate the error caused by replacing $g_j$ with $h_j$. Using
$$|x^m-y^m|\le m|x-y|\max(|x|,|y|)^{m-1}, $$
we obtain
\begin{equation*}
\begin{aligned}
&\left\|\prod_{j=1}^2 \Bigl(S_N(\bar g_j)\Bigr)^{m_j}-\prod_{j=1}^2 \Bigl(S_N(\bar h_j)\Bigr)^{m_j}\right\|_{L^1}\\
&\le m_1\left\|S_N(\bar g_1)-S_N(\bar h_1)\right\|_{L^1}\max\left(\left\|S_N(\bar g_1)\right\|_{L^\infty},
\left\|S_N(\bar h_1)\right\|_{L^\infty}\right)^{m_1-1}\left\|S_N(\bar g_2)\right\|_{L^\infty}^{m_2} \\
&\quad+m_2\left\|S_N(\bar g_2)-S_N(\bar h_2)\right\|_{L^1}\left\|S_N(\bar h_1)\right\|_{L^\infty}^{m_1}
\max\left(\left\|S_N(\bar g_2)\right\|_{L^\infty},\left\|S_N(\bar h_2)\right\|_{L^\infty}\right)^{m_2-1}.
\end{aligned}
\end{equation*}
Using \eqref{cL1e},
$$\left\|S_N(\bar g_1)-S_N(\bar h_1)\right\|_{L^1}\ll N\epsilon_N r_N^a,\quad
\left\|S_N(\bar g_2)-S_N(\bar h_2)\right\|_{L^1}\ll N\epsilon_N r_N^{a-\beta}.$$
Moreover,
$$\left\|S_N(\bar g_1)\right\|_{L^\infty},\left\|S_N(\bar h_1)\right\|_{L^\infty}\ll N,
\quad\text{and}\quad
\left\|S_N(\bar g_2)\right\|_{L^\infty},\left\|S_N(\bar h_2)\right\|_{L^\infty}\ll Nr_N^{-\beta}.$$
Therefore
$$\left\|\prod_{j=1}^2 \Bigl(S_N(\bar g_j)\Bigr)^{m_j}-\prod_{j=1}^2 \Bigl(S_N(\bar h_j)\Bigr)^{m_j}\right\|_{L^1}\ll \epsilon_N N^m r_N^{a-m_2\beta}.
$$
After division by the normalization in \eqref{momentclaim}, this is bounded by
$O\bigl(\epsilon_N N^{\frac{m}{2}}r_N^{a-\frac{md}{2}}\bigr).$
Since $Nr_N^d\to\infty$, choosing $\epsilon_N=N^{-B}$ with $B>0$ sufficiently large makes this term $o(1)$. Thus the error caused by replacing  $g_j$ by $h_j$ in \eqref{momentclaim} is negligible.




Therefore, the claim \eqref{momentclaim} will follow once we show
\begin{equation}\label{hbarerr}
\lim_{N\rightarrow\infty} \frac{\int_M \left(S_N(\bar{h}_1) \right)^{m_1} \left(S_N(\bar{h}_2)\right)^{m_2}\d\mu}{\left(B_d \rho(x^*)\right)^{\frac{m}{2}} N^{\frac{m}{2}} r_N^{\frac{md}{2}-m_2\beta }}=\omega(m_1,m_2).
\end{equation}


Step 2: Cluster decomposition.

We expand the product
\begin{align*}
\prod_{j=1}^2 (S_N(\bar{h}_j))^{m_j} &= \sum_{\textbf{n}\in \{0,\ldots,N-1\}^m}
\left(\prod_{i=1}^{m_1} (\bar{h}_1 \circ T^{n_i})\right) \left(\prod_{i=m_1+1}^m (\bar{h}_2 \circ T^{n_i})\right)\\
&=\sum_{\textbf{n}\in \{0,\ldots,N-1\}^m} \prod_{i=1}^{m} (\bar{h}_{j(i)} \circ T^{n_i}),
\end{align*}
where
\begin{equation*}
j(i)=\begin{cases}
1 & \text{if } 1\leq i \leq m_1,\\
2 & \text{if } m_1 + 1 \leq i \leq m.
\end{cases}
\end{equation*}



Let $ C > 0 $ denote a sufficiently large constant. For each $\textbf{n}= (n_1, \dots, n_m)\in \{0,\ldots,N-1\}^m$, define clusters as follows. There exists a minimal integer $1 \leq k=k(\textbf{n}) \leq m^2 + 1$ such that for all $i, i'\in \{1,\dots,m\},$
\begin{equation}\label{nigp}
|n_i - n_{i'}| \leq C^k \log N \quad \text{or} \quad |n_i - n_{i'}| > C^{k+1} \log N.
\end{equation}
For $l\geq 1$ we call $\{i_1,\ldots,i_l\}\subset \{1,\ldots,m\}$ a cluster of $\textbf{n}$ if it is a maximal collection with
\begin{equation*}
|n_i - n_{i'}| < C^k \log N \quad \forall i,i'\in \{i_1,\ldots,i_l\}.
\end{equation*}
In such cases, define $ l $ as the cluster length. Then there exists an integer $1\leq \eta = \eta(\textbf{n}) \leq m$ such that $ \textbf{n}$ consists of the clusters $\{i^s_1, \dots, i^s_{l_s}\}$ for $ s = 1, \dots, \eta $, with corresponding lengths $ l_s $.


Equivalently, let $ k $ be the smallest integer for which the set $\{n_i\}$ can be partitioned into $\eta$ clusters of sizes $l_s$, each being a maximal subset where pairwise differences are at most $C^k \log N$, and distinct clusters are separated by more than $C^{k+1} \log N$.

By exponential mixing, we have
\begin{equation}\label{cluexp}
\begin{aligned}
\left|\int \left(\prod_{i=1}^m (\bar{h}_{j(i)} \circ T^{n_i})\right) \d\mu \right.&\left.-\prod_{s=1}^{\eta}\int\left(\prod_{r=1}^{l_s} (\bar{h}_{j(i^s_r)} \circ T^{n_{i^s_r}})\right) \d\mu\right|\\
&\ll e^{-\gamma C^{k+1} \log N}\prod_{s=1}^\eta\left\|\prod_{r=1}^{l_s} (\bar{h}_{j(i^s_r)} \circ T^{n_{i^s_r}})\right\|_{C^{\kappa}}.
\end{aligned}
\end{equation}
Each $\bar{h}_j$ satisfies $\|\bar{h}_j\|_{C^{\kappa}} \ll N^{B'}$ for a large constant $B'>0$. Since the time differences within each cluster are at most $C^k \log N$, we obtain
\begin{align*}
\left\|\prod_{r=1}^{l_s} (\bar{h}_{j(i^s_r)} \circ T^{n_{i^s_r}})\right\|_{C^{\kappa}} &\ll\prod_{r=1}^{l_s} \| \bar{h}_{j(i^s_r)} \|_{C^{\kappa}} L^{C^k \log N}&\ll\left(N^{B'} L^{C^k \log N}\right)^{l_s},
\end{align*}
where $L=\sup_{\|f\|_{C^{\kappa}}=1} \|f\circ T\|_{C^{\kappa}}$ is the operator norm of $f\mapsto f\circ T$ on $C^{\kappa}$.

Consequently, the right-hand side of \eqref{cluexp} can be estimated as
\begin{equation}\label{proderr}
e^{-\gamma C^{k+1}\log N}N^{B'm}L^{mC^k\log N}\le N^{-100m},
\end{equation}
which holds for sufficiently large $C$ independent of $N.$ After summing over all $\mathbf n$, this error is negligible.




Considering \eqref{hbarerr} and \eqref{proderr}, the claim \eqref{momentclaim} will follow once we show
\begin{equation*}
\lim_{N\rightarrow \infty} \frac{\sum_{\textbf{n}\in \{0,\ldots,N-1\}^m} \prod_{s=1}^{\eta(\textbf{n})} \int_M \prod_{r=1}^{l_s} (\bar{h}_{j(i^s_r)} \circ T^{n_{i^s_r}}) \d\mu}{\left(B_d \rho(x^*)\right)^{\frac{m}{2}} N^{\frac{m}{2}} r_N^{\frac{md}{2}-m_2\beta }} = \omega(m_1,m_2).
\end{equation*}

Replacing once more the smooth approximations $\bar{h}_{j(i^s_r)}$ by $\bar{g}_{j(i^s_r)}$ on the left hand side, similarly to \eqref{hbarerr}, only introduces another negligible error term of order $o(1)$. Therefore, the proof will be finished once we show
\begin{equation}\label{momentsumclaim}
\lim_{N\rightarrow \infty} \frac{\sum_{\textbf{n}\in \{0,\ldots,N-1\}^m} \prod_{s=1}^{\eta(\textbf{n})} \int_M \prod_{r=1}^{l_s} (\bar{g}_{j(i^s_r)} \circ T^{n_{i^s_r}})\d\mu}{\left(B_d \rho(x^*)\right)^{\frac{m}{2}} N^{\frac{m}{2}} r_N^{\frac{md}{2}-m_2\beta }} = \omega(m_1,m_2).
\end{equation}

Step 3: Negligible cluster configurations.

Define $\Delta_\eta = \left\{\mathbf{n} : \mathbf{n} \text{ consists of } \eta \text{ clusters}\right\}.$
It holds that
\begin{equation}\label{clusternumber}
\#\Delta_{\eta} \ll \binom{m}{\eta}\eta !N^{\eta}\eta^{m-\eta} \left(2C^k\log N\right)^{m-\eta}\ll N^{\eta} (\log N)^{m}.
\end{equation}



Case 3.1: $\eta>\frac{m}{2}.$ At least one cluster is a singleton.
Since
$$\int\bar{g}_{j(i)} \circ T^{n_i} \d\mu = 0,$$
then it holds that
\begin{equation*}
\prod_{s=1}^{\eta} \int_M \prod_{r=1}^{l_s} (\bar{g}_{j(i^s_r)} \circ T^{n_{i^s_r}}) \d\mu=0.
\end{equation*}

Case 3.2: $\eta<\frac{m}{2}.$ If at least one cluster is a singleton, it can be treated as in Case 3.1.
If no cluster is a singleton, then all $l_s \geq 2$. Given that
\begin{align*}
&\|\bar{g}_1\|_{L^2} \ll r_N^{\frac{d}{2}},\quad \|\bar{g}_1\|_{L^{\infty}} \ll 1, \quad \text{ and}\\
&\|\bar{g}_2\|_{L^2} \ll r_N^{-\beta+\frac{d}{2}},\quad \|\bar{g}_2\|_{L^{\infty}} \ll r_N^{-\beta},
\end{align*}




we have
\begin{equation}\label{clusterint}
\left|\prod_{s=1}^{\eta}\int\prod_{r=1}^{l_s}(\bar{g}_{j(i^s_r)} \circ T^{n_{i^s_r}})\d\mu\right|\ll
r_N^{-m_2 \beta + \eta d}.
\end{equation}
Thus
\begin{equation*}
\begin{aligned}
\sum_{\eta<m/2} \sum_{\mathbf{n} \in \Delta_\eta}
\left|\prod_{s=1}^{\eta}\int\prod_{r=1}^{l_s}(\bar{g}_{j(i^s_r)} \circ T^{n_{i^s_r}})\d\mu\right|
&\ll\sum_{\eta<m/2}N^{\eta} (\log N)^m r_N^{-m_2 \beta + \eta d}\\
&\ll\sum_{\eta<m/2} \left(N r_N^d\right)^\eta r_N^{- m_2 \beta}(\log N)^m.
\end{aligned}
\end{equation*}
Subcase 3.2~(i): $N r_N^d>N^{1/100}$. This bound is of order $o\left(N^{\frac{m}{2}} r_N^{\frac{md}{2} - m_2 \beta}\right).$

Subcase 3.2~(ii): $N r_N^d\leq N^{1/100}$. For each index $\textbf{n}\in \Delta_{\eta}$ we distinguish two cases:

Subcase 3.2~(ii-a): all $n_i$ within the same cluster are equal for each cluster.
Then there exist $O(N^\eta)$ such indices \textbf{n}. Taking into account the estimate in \eqref{clusterint}, the total contribution from these indices is of order $O\left(N^{\eta} r_N^{\eta d - m_2 \beta}\right).$



Subcase 3.2~(ii-b): at least one cluster contains two distinct indices. Then Lemma \ref{slowrecanlem} provides an additional saving. Using the estimates established in Case~4.1~(B) below,
\eqref{clusterint} can be replaced by
\begin{align*}
\left| \prod_{s=1}^{\eta} \int \prod_{r=1}^{l_s} (\bar{g}_{j(i^s_r)} \circ T^{n_{i^s_r}} )\d\mu \right| &\ll
r_N^{\eta d - m_2 \beta} |\log r_N|^{-100m} \\
&\ll
r_N^{\eta d - m_2 \beta} (\log N)^{-100m}.
\end{align*}
Note that here the assumption $N r_N^d\leq N^{\frac{1}{100}}$ is used crucially to show $\log N \asymp |\log r_N|$.

Since the number of indices $\mathbf{n}$ in this subcase is at most $\#\Delta_{\eta}\ll N^{\eta} (\log N)^m$, the total contribution from this case is of order
\begin{equation*}
O\left(N^{\eta} r_N^{\eta d - m_2 \beta} (\log N)^{-99m}\right).
\end{equation*}


Combining subcases 3.2(ii-a) and 3.2(ii-b), it holds that
\begin{align*}
\sum_{\eta < \frac{m}{2}} \sum_{\textbf{n} \in \Delta_{\eta}} \left| \prod_{s=1}^{\eta} \int \prod_{r=1}^{l_s} (\bar{g}_{j(i^s_r)} \circ T^{n_{i^s_r}}) \d\mu\right| &\ll \sum_{\eta < \frac{m}{2}} N^{\eta} r_N^{\eta d - m_2 \beta}.
\end{align*}
Using the assumption $N r_N^d \rightarrow\infty$, it follows that $N^{\eta} r_N^{\eta d - m_2 \beta} =o\left( N^{\frac{m}{2}} r_N^{\frac{md}{2} - m_2 \beta}\right)$ for each $\eta<\frac{m}{2}$. Therefore the above bound is of order $o\left( N^{\frac{m}{2}} r_N^{\frac{md}{2} - m_2 \beta}\right)$.


In conclusion, all terms in \eqref{momentsumclaim} for $\eta \neq \frac{m}{2}$ are asymptotically negligible. Note that, in particular, if $m$ is odd, then \eqref{momentsumclaim} is satisfied. So for the rest of the proof, assume $m$ is even.

Step 4: The critical case $\eta=\frac{m}{2}$.

Unless all the clusters are of length $2$, $\textbf{n}$ must have a singleton. Define
$$
\Gamma=\left\{\mathbf{n}\in \Delta_{\frac{m}{2}} : \text{each cluster has size } 2\right\}.
$$
Then we consider
$\sum_{\mathbf{n}\in\Gamma}\prod_{s=1}^{\frac{m}{2}}\int \big(\bar{g}_{j(i^s_1)} \cdot \bar{g}_{j(i^s_2)} \circ T^{|n_{i^s_2}-n_{i^s_1}|}\big)\d\mu.$

Since all terms with singletons are negligible -- as in Case 3.1 -- it is enough to show
\begin{equation}\label{momentgammaclaim}
\lim_{N\rightarrow \infty} \frac{\sum_{\textbf{n}\in \Gamma} \prod_{s=1}^{\frac{m}{2}} \int_M \prod_{r=1}^{2} (\bar{g}_{j(i^s_r)} \circ T^{n_{i^s_r}}) \d\mu}{\left(B_d \rho(x^*)\right)^{\frac{m}{2}} N^{\frac{m}{2}} r_N^{\frac{md}{2}-m_2\beta }} = \omega(m_1,m_2).
\end{equation}

Step 4.1: Estimates for a single pair.

For $n_i,n_{i'}$ in the same cluster, we aim to estimate $\int_M (\bar{g}_{j(i)} \circ T^{n_i})(\bar{g}_{j(i')} \circ T^{n_{i'}})\d\mu$. Simplifying notation, we consider
\begin{equation*}
\int_M \bar{g}_j \cdot \bar{g}_{j'} \circ T^k\d\mu,
\end{equation*}
for $j,j' \in \{1,2\}$ and $k=|n_i - n_{i'}| \leq C^{m^2+1} \log N,$ see \eqref{nigp}. Below we will distinguish several cases.


We first highlight a technical point. Since $g_1 = \one_{A_N}$, the slow recurrence estimate, specifically Lemma~\ref{slowrecanlem}, is applicable only up to times $k=C' |\log r_N|$ for a constant $C'>0$, which might be small if $r_N$ decays substantially slower than $1/N$.






Case 4.1~(A): $k=0$. We consider the cases where $j=j'=1,$ $j=j'=2$, and $j\neq j'$ respectively. For $j=j'=1$ we use Lemma~\ref{regadaptmeasure} to obtain
\begin{equation}\label{mainest1}
\int \bar{g}_1^2\d\mu\sim \mu(A_N)\sim B_d \rho(x^*) r_N^d.
\end{equation}
For $j=j'=2$, a straightforward calculation in Euclidean space reveals
\begin{equation}\label{mainest2}
\begin{aligned}
\int_M \bar{g}_2^2 \d\mu &\sim\int_{r_N \leq d(x,x^*) < 1} f^2(x)\d\mu(x)\\
&\sim \rho(x^*)\int_{\textbf{y}\in \R^d, r_N < \|\textbf{y}\| < 1} \frac{1}{\|\textbf{y}\|^{2\beta}} \d \textbf{y}\\
&\sim B_d \rho(x^*) r_N^{-2\beta+d} \int_1^{\infty} u^{-2\alpha} \d u \\
&\sim \sigma^2 B_d \rho(x^*) r_N^{-2\beta+d},
\end{aligned}
\end{equation}
where $\sigma^2=\frac{1}{2\alpha-1}$.


A straightforward computation in local coordinates near $x^*$, together with the boundedness of $f$ away from $x^*$, gives
\begin{equation*}
\|g_2\|_{L^1}\ll \begin{cases}
1 & \text{if } \beta<d\\
|\log r_N| & \text{if } \beta=d\\
r_N^{d-\beta} & \text{if } \beta>d.
\end{cases}
\end{equation*}
For $j\neq j'$, since the supports of $g_1$ and $g_2$ are disjoint, we therefore have
\begin{align*}
\left|\int_M \bar g_1\bar g_2\d\mu\right|=\left|\int_M g_1\d\mu\right|\left|\int_M g_2\d\mu\right|
&\ll r_N^{d}\left(1+|\log r_N|+r_N^{d-\beta}\right)\\
&=r_N^{d-\beta} \left(r_N^{\beta}(1+|\log r_N|) + r_N^d\right).
\end{align*}






Case 4.1~(B): $0<k<C' \left|\log r_N \right|,$ where $C'>0$ is a sufficiently large constant. For $j=j'=1$, applying Lemma~\ref{slowrecanlem} to $A_N \cap T^{-k}(A_N)$ with $A=100m$ yields
\begin{equation*}
\left|\int_M \bar{g}_1 \cdot \bar{g}_1 \circ T^{k} \d\mu\right|\ll \mu\left(A_N \cap T^{-k}(A_N)\right)
\ll r_N^d \left|\log r_N \right|^{-100m}.
\end{equation*}

For $j=j'=2$, Lemma \ref{truncvarlem} with $A=100m$ yields
\begin{equation}
\left| \int_M \bar{g}_2 \cdot \bar{g}_2 \circ T^k \d\mu \right| \ll r_N^{-2\beta+d} |\log r_N|^{-100m}.
\end{equation}

For the cross term $j\neq j'$, we consider the case $j=1$ and $j'=2$;
the reverse situation follows analogously.

Let $\hat{r}_N=r_N|\log r_N|^H$, where $H>\frac{100m}{\beta}$ is fixed. Since $A_N$ is regularly adapted, for sufficiently large $N$ we have $A_N\subset B_{\hat{r}_N}(x^*)$. Hence,
\begin{equation*}
\int_M g_1\cdot g_2\circ T^k\d\mu
=
\int_{A_N\cap T^{-k}(M\setminus A_N)}f\circ T^k\d\mu
=
\mathfrak A+\mathfrak B,
\end{equation*}
where
\begin{align*}
\mathfrak A&:=
\int_{A_N\cap T^{-k}(B_{\hat{r}_N}(x^*)\setminus A_N)}
f\circ T^k\d\mu,\\
\mathfrak B&:=
\int_{A_N\cap T^{-k}(M\setminus B_{\hat{r}_N}(x^*))}
f\circ T^k\d\mu .
\end{align*}

For $\mathfrak A$, note that for sufficiently large $N$, $|\log r_N|\leq 2|\log\hat{r}_N|$. Since $k<C'|\log r_N|$, we therefore have $k<2C'|\log\hat r_N|.$
Using slow recurrence with exponent $A'=Hd+100m$ and the inclusion $B_{r_N}(x^*)\subset A_N\subset B_{\hat r_N}(x^*),$
we obtain
\begin{align*}
\mathfrak A
&\leq
\sup_{y\notin B_{r_N}(x^*)}|f(y)|
\mu(B_{\hat{r}_N}(x^*)\cap T^{-k}B_{\hat{r}_N}(x^*))\\
&\ll
r_N^{-\beta}\hat{r}_N^d|\log\hat{r}_N|^{-A'}\\
&\ll
r_N^{d-\beta}|\log r_N|^{-100m}.
\end{align*}

For $\mathfrak B$, using Lemma~\ref{regadaptmeasure},
$$\mathfrak B\leq\sup_{y\notin B_{\hat{r}_N}(x^*)}|f(y)|\,\mu(A_N)\ll
\hat{r}_N^{-\beta}r_N^d\ll r_N^{d-\beta}|\log r_N|^{-100m}.
$$

Combining the estimates for $\mathfrak A$ and $\mathfrak B$ gives
\begin{equation*}
\int_M g_1\cdot g_2\circ T^k\d\mu
\ll
r_N^{d-\beta}|\log r_N|^{-100m}.
\end{equation*}
Finally, centering does not affect the estimate, so for $j\neq j'$,
\begin{equation*}
\left|\int_M\bar g_j\cdot\bar g_{j'}\circ T^k\d\mu\right|
\ll
r_N^{d-\beta}|\log r_N|^{-100m}.
\end{equation*}


Case 4.1~(C): $k> C' \left|\log r_N\right|.$ We emphasize that this case only needs to be considered when $|\log r_N|=o(\log N)$. Indeed, all computations here are still inside a cluster, so the relevant time differences satisfy $k=O(\log N)$. If $|\log r_N|$ is comparable to $\log N$, then taking $C'$ large enough makes the condition $k>C'|\log r_N|$ incompatible with the intra-cluster bound $k=O(\log N)$, and this subcase becomes empty. Using \eqref{regasm2} with some $\varepsilon\in (0,r_N)$ to be determined later, there exists a $C^{\kappa}$ smooth function $\theta:M\rightarrow[0,1]$ such that
\begin{equation*}
\|g_1-\theta\|_{L^1} \ll\varepsilon r_N^a \quad \text{and} \quad \|\theta\|_{C^{\kappa}} \ll\varepsilon^{-\kappa} r_N^{-a'}.
\end{equation*}
Denote $\bar{\theta}=\theta - \int_M \theta\d\mu$. By applying exponential mixing, we obtain
\begin{align*}
\left| \int_M \bar{g}_{1} \cdot \bar{g}_{1} \circ T^{k} \d\mu \right| &\ll \left| \int_M \bar{\theta} \cdot \bar{\theta} \circ T^{k} \d\mu \right| + \varepsilon r_N^{a} \\
& \ll e^{-\gamma k} \varepsilon^{-2\kappa} r_N^{-2a'} +\varepsilon r_N^a.
\end{align*}
Choosing $\varepsilon=e^{-\frac{\gamma k}{2\kappa+1}}$ guarantees $\varepsilon<r_N$ whenever $C'>(2\kappa+1)/\gamma$, since $k>C'\left|\log r_N\right|$. We then obtain
\begin{equation*}
\left| \int_M \bar{g}_{1} \cdot \bar{g}_{1} \circ T^{k} \d\mu \right| \ll e^{-\frac{\gamma k}{2\kappa+1}} r_N^{-2a'}.
\end{equation*}
If $k=C'|\log r_N| + i$ for $i\in\mathbb{N},$ then we have
\begin{equation*}
\left| \int_M \bar{g}_{1} \cdot \bar{g}_{1} \circ T^{k} \d\mu \right| \ll r_N^{\frac{\gamma C'}{2\kappa+1}-2a'}e^{-\frac{\gamma i}{2\kappa+1}}\ll r_N^{100d} e^{-\frac{\gamma i}{2\kappa+1}}
\end{equation*}
for large $C'$.


The same argument, with the corresponding smooth approximations of $g_2$ and of the cross terms, gives
\begin{equation*}
\left| \int_M \bar{g}_{2} \cdot \bar{g}_{2} \circ T^{k} \d\mu \right| \ll r_N^{100d} e^{-\frac{\gamma i}{2\kappa+1}}
\end{equation*}
and
\begin{equation*}
\left| \int_M \bar{g}_{j} \cdot \bar{g}_{j'} \circ T^{k} \d\mu \right| \ll r_N^{100d} e^{-\frac{\gamma i}{2\kappa+1}},
\end{equation*}
for $j\neq j'$.







Step 4.2: Negligibility outside the diagonal pair configurations.

Let
\begin{equation*}
\Gamma'=\left\{\textbf{n} \in \Gamma \;|\; n_{i^s_1}=n_{i^s_2} \text{ and } j(i^s_1)=j(i^s_2) \;\; \forall s=1,\ldots,\frac{m}{2}\right\},
\end{equation*}
then, using the estimates from $j\neq j'$ in case 4.1~(A), and cases 4.1~(B) and 4.1~(C) above, we obtain
\begin{align*}
&\sum_{\textbf{n} \in \Gamma\setminus \Gamma'} \prod_{s=1}^{\frac{m}{2}} \left|\int_M \bar{g}_{j(i^s_1)} \cdot \bar{g}_{j(i^s_2)} \circ T^{n_{i^s_2}-n_{i^s_1}} \d\mu\right|\\
&\ll N^{\frac{m}{2}} r_N^{\frac{md}{2}-m_2 \beta}\left(r_N^\beta(1+|\log r_N|)+r_N^d \right)
+ N^{\frac{m}{2}} \left|\log r_N \right|^{\frac{m}{2}} r_N^{\frac{md}{2} - m_2 \beta} \left|\log r_N \right|^{-100m}\\
&\ll N^{\frac{m}{2}} r_N^{\frac{md}{2}-m_2 \beta} \left|\log r_N \right|^{-99m}.
\end{align*}
Note that here we used the fact that all cases, where one of the $|n_{i^s_1}-n_{i^s_2}|>C'|\log r_N|$, are negligible by summing the estimate for $i$ from Case 4.1~(C).




Step 4.3: The diagonal pair contribution.

Using \eqref{mainest1} and \eqref{mainest2}, for $\mathbf{n}\in \Gamma'$ it holds that
$$
\prod_{s=1}^{\frac{m}{2}} \int_M \bar{g}_{j(i^s_1)}^2 \d\mu\sim\left(B_d \rho(x^*)\right)^{\frac{m}{2}} r_N^{\frac{md}{2}-m_2\beta } \sigma^{m_2}.
$$
If either $m_1$ or $m_2$ is odd, then $\Gamma'=\varnothing$, since indices of type $g_1$ and $g_2$ must be paired separately. Otherwise, there are $(m_1-1)!!$ and $(m_2-1)!!$ possible pairings within each type. For each fixed pairing, this leaves $m/2$ free time variables, hence $N^{\frac{m}{2}}$ choices. The cluster-separation condition excludes only configurations where two such times lie within distance $O(\log N)$, whose number is $O(N^{\frac{m}{2}-1}\log N)$. Hence
$$\#\Gamma'\sim N^{\frac{m}{2}}(m_1-1)!!(m_2-1)!!.
$$
It follows that
$$
\sum_{\mathbf{n} \in \Gamma'} \prod_{s=1}^{\frac{m}{2}} \int_M \bar{g}_{j(i^s_1)}^2 \d\mu \sim \left(B_d \rho(x^*)\right)^{\frac{m}{2}} N^{\frac{m}{2}} r_N^{\frac{md}{2}-m_2\beta } \omega(m_1,m_2),
$$
which completes the proof of \eqref{momentgammaclaim}.
\end{proof}

\end{document}
